\documentclass{article}
\usepackage{../assignment_preamble}

\title{Homework 2}
\author{Ravi Kini}
\date{February 18, 2026}

\begin{document}

\maketitle

\problem{}
\subproblem{(a)}
Let $A$ be a symmetric $3 \times 3$ matrix:
\begin{align*}
    A = \begin{bmatrix}
        a_1 & b & c \\ b & a_2 & d \\ c & d & a_3
    \end{bmatrix}
\end{align*}
To annihilate the $c$ entries (in row/column $3$ and column/row $1$), we use a Givens rotation that acts as the identity except in rows and columns $2, 3$. Letting $r = \sqrt{b^2 + c^2}$, the rotation matrix $Q$ has the form:
\begin{align*}
    Q & = \begin{bmatrix}
        1 & 0 & 0 \\ 0 & b/r & c/r \\ 0 & -c/r & b/r
    \end{bmatrix}
\end{align*}
Then, to see that it annihilates the $c$ entries:
\begin{align*}
    (QAQ_T)_{1,3} & = \sum_{i=1}^3 (QA)_{1,i}Q^T_{i,3} = \sum_{i=1}^3 \sum_{j=1}^3 Q_{1,j}A_{j,i}Q^T_{i,3} \\
    & = \sum_{i=1}^3 Q_{1,1}A_{1,i}Q^T_{i,3} = \sum_{i=2}^3 Q_{1,1}A_{1,i}Q^T_{i,3} \\
    & = 1 \cdot b \cdot -c/r + 1 \cdot c \cdot b/r = 0 \\
    (QAQ_T)_{3,1} & = \sum_{i=1}^3 (QA)_{3,i}Q^T_{i,1} = \sum_{i=1}^3 \sum_{j=1}^3 Q_{3,j}A_{j,i}Q^T_{i,1} \\
    & = \sum_{i=1}^3 \sum_{j=2}^3 Q_{3,j}A_{j,i}Q^T_{i,1} = \sum_{j=2}^3 Q_{3,j}A_{j,1}Q^T_{1,1} \\
    & = -c/r \cdot b \cdot 1 + b/r \cdot c \cdot 1 = 0 \\
\end{align*}
\subproblem{(b)}
Let $A$ be a symmetric $n \times n$ matrix where the leading $(n - 1) \times (n - 1)$ submatrix has been tridiagonalized:
\begin{align*}
    A = \begin{bmatrix}
        a_1 & b_2 & & & & \\
        b_2 & a_2 & b_3 & & & \\
        & b_3 & a_3 & \ddots & & \\
        & & \ddots & \ddots & b_{n-1} & d \\
        & & & b_{n-1} & a_{n-1} & c \\
        & & & d & c & a_n
    \end{bmatrix}
\end{align*}
To annihilate the $d$ entries (in row/column $n$ and column/row $n - 2$), we use a Givens rotation that acts as the identity except in rows and columns $n - 1, n$. Letting $r = \sqrt{b_{n-1}^2 + d^2}$, the rotation matrix $Q$ has the form:
\begin{align*}
    Q & = \begin{bmatrix}
        I_{n-2} & 0 & 0 \\ 0 & b_{n-1}/r & d/r \\ 0 & -d/r & b_{n-1}/r
    \end{bmatrix}
\end{align*}
We see it is orthogonal, and that the product $QAQ^T$ would have the same leading $(n - 2) \times (n - 2)$ submatrix. Then, to see that it annihilates the $d$ entries:
\begin{align*}
    (QAQ_T)_{n-2,n} & = \sum_{i=1}^n (QA)_{n-2,i}Q^T_{i,n} = \sum_{i=1}^n \sum_{j=1}^n Q_{n-2,j}A_{j,i}Q^T_{i,n} \\
    & = \sum_{i=1}^n Q_{n-2,n-2}A_{n-2,i}Q^T_{i,n} = \sum_{i=n-1}^n Q_{n-2,n-2}A_{n-2,i}Q^T_{i,n} \\
    & = 1 \cdot b_{n-1} \cdot -d/r + 1 \cdot d \cdot b_{n-1}/r = 0 \\
    (QAQ_T)_{n,n-2} & = \sum_{i=1}^n (QA)_{n,i}Q^T_{i,n-2} = \sum_{i=1}^n \sum_{j=1}^n Q_{n,j}A_{j,i}Q^T_{i,n-2} \\
    & = \sum_{i=1}^n \sum_{j=n-1}^n Q_{n,j}A_{j,i}Q^T_{i,n-2} = \sum_{j=n-1}^n Q_{n,j}A_{j,n-2}Q^T_{n-2,n-2} \\
    & = -d/r \cdot b_{n-1} \cdot 1 + b_{n-1}/r \cdot d \cdot 1 = 0 \\
\end{align*}

\clearpage
\problem{}
Let $T$ be a symmetric tridiagonal $n \times n$ matrix, with QR factorization $T = QR$. Consider the matrix $RQ$. We see that:
\begin{align*}
    R & = Q^T T \\
    RQ & = Q^T T Q \\
    (RQ)^T & = (Q^T T Q)^T \\
    & = Q^T T^T Q = Q^T T Q = RQ
\end{align*}
Since $(RQ)^T = RQ$, we see that $RQ$ is symmetric.

Now note that in QR factorization, column $k$ of $Q$ is column $k$ of $A$ orthogonalized with respect to the previous columns, and so can only have nonzero components in the span of those columns (and itself). Since the first $k$ columns of $A$ have zeros in rows $i$ where $i > k + 1$, it follows that column $k$ of $Q$ must also have a zero in such rows. Equivalently, $Q_{i,j} = 0$ for $i > j + 1$. In other words, $Q$ is zero below the first subdiagonal. Further, since $R$ is upper triangular, it is zero below the diagonal ($R_{i,j} = 0$ for $i > j$). We now consider $(RQ)_{i,j}$, with $i > j + 1$:
\begin{align*}
    (RQ)_{i,j} & = \sum_{k=1}^n R_{i,k}Q_{k,j} \\
    & = \sum_{k=i}^n R_{i,k}Q_{k,j} \\
    & = \sum_{k=i}^n R_{i,k} \cdot 0 = 0
\end{align*}
Since $(RQ)_{i,j} = 0$ for $i > j + 1$, it follows that $RQ$ is zero below the first subdiagonal. Since $RQ$ is symmetric, it must alo be zero above the first superdiagonal, which means that $RQ$ is tridiagonal symmetric.

\clearpage
\problem{}
Consider the shifted QR iteration algorithm. Let $\underline{Q_k} = Q_1Q_2 \ldots Q_k$ and $\underline{R_k} = R_kR_{k-1} \ldots R_1$. We proceed by induction. For the base case where $k = 1$, we see that:
\begin{align*}
    Q_1R_1 & = A - \mu_1 I \\
    A_1 & = R_1Q_1 + \mu_1 I \\
    & = Q_1^TQ_1R_1Q_1 + \mu_1 I \\
    & = Q_1^T(A - \mu_1 I)Q_1 + \mu_1 I \\
    & = Q_1^TAQ_1 - \mu_1 I + \mu_1 I = Q_1AQ_1
\end{align*}
Now assume that for some $k$, $A_k = \underline{Q_k}^TA\underline{Q_k}$ and $(A - \mu_k I) \ldots (A - \mu_1 I) = \underline{Q_k}\underline{R_k}$. Then:
\begin{align*}
    Q_{k+1}R_{k+1} & = A_k - \mu_{k+1}I \\
    \underline{Q_{k+1}}\underline{R_{k+1}} & = \underline{Q_k}Q_{k+1}R_{k+1}\underline{R_k} \\
    & = \underline{Q_k}(A_k - \mu_{k+1}I)\underline{R_k} \\
    & = \underline{Q_k}(\underline{Q_k}^TA\underline{Q_k} - \mu_{k+1}I)\underline{R_k} = (A\underline{Q_k} - \mu_{k+1}\underline{Q_k}I)\underline{R_k} \\
    & = (A - \mu_{k+1}I)\underline{Q_k}\underline{R_k} \\
    & = (A - \mu_{k+1} I)(A - \mu_k I) \ldots (A - \mu_1 I)
\end{align*}
Further:
\begin{align*}
    A_{k+1} & = R_{k+1}Q_{k+1} + \mu_{k+1} I \\
    & = Q_{k+1}^TQ_{k+1}R_{k+1}Q_{k+1} + \mu_{k+1} I \\
    & = Q_{k+1}^T(A_k - \mu_{k+1}I)Q_{k+1} + \mu_{k+1} I \\
    & = Q_{k+1}^TA_kQ_{k+1} - \mu_{k+1}I + \mu_{k+1} I \\
    & = Q_{k+1}^T\underline{Q_k}^TA\underline{Q_k}Q_{k+1} \\
    & = \underline{Q_{k+1}}^TA\underline{Q_{k+1}}
\end{align*}
By induction, it follows that $A_k = \underline{Q_k}^TA\underline{Q_k}$ and $(A - \mu_k I) \ldots (A - \mu_1 I) = \underline{Q_k}\underline{R_k}$ for all $k$.

\clearpage
\problem{}
\subproblem{(a)}
\begin{lstlisting}[language=Python]
import math
EPS = 1e-15

def eye(n):
    return [[float(i == j) for i in range(n)] for j in range(n)]

def copy(A):
    return [row[:] for row in A]

def givens_left(A, i, j, c, s):
    n = len(A)
    for k in range(n):
        a, b = A[i][k], A[j][k]
        A[i][k] = c * a + s * b
        A[j][k] = -s * a + c * b

def givens_right(A, i, j, c, s):
    n = len(A)
    for k in range(n):
        a, b = A[k][i], A[k][j]
        A[k][i] = c * a - s * b
        A[k][j] = s * a + c * b

def tridiagonalize(A):
    n = len(A)
    T = copy(A)
    for i in range(n - 2):
        for j in range(n - 1, i + 1, -1):
            b = T[i + 1][i]
            c = T[j][i]
            if abs(c) < EPS:
                continue
            r = math.hypot(b, c)
            givens_left(T, i + 1, j, b / r, c / r)
            givens_right(T, i + 1, j, b / r, -c / r)
    return T
\end{lstlisting}
\subproblem{(b)}
\begin{lstlisting}
def QR(T):
    n = len(T)
    Q = eye(n)
    R = copy(T)
    for k in range(n - 1):
        a = R[k][k]
        b = R[k + 1][k]
        if abs(b) < EPS:
            continue
        r = math.hypot(a, b)
        givens_left(R, k, k + 1, a / r, b / r)
        givens_right(Q, k, k + 1, a / r, -b / r)
    return Q, R
\end{lstlisting}
\subproblem{(c)}
\begin{lstlisting}
def leading(A):
    n = len(A)
    return [row[:-1] for row in A[:-1]]

def zeros(n):
    return [[0.0 for i in range(n)] for j in range(n)]
    
def wilkinson(T):
    n = len(T)
    a, b, c = T[n - 2][n - 2], T[n - 2][n - 1], T[n - 1][n - 1]
    delta = (a - c) / 2
    sign = 1.0 if delta >= 0 else -1.0
    return c - sign * b ** 2 / (abs(delta) + math.hypot(delta, b))

def QR_iteration(T, tol=1e-10, max_iter=1000):
    n = len(T)
    T = copy(T)
    eigs = []
    while n > 1:
        for _ in range(max_iter):
            if abs(T[n - 1][n - 2] / (abs(T[n - 2][n - 2]) + abs(T[n - 1][n - 1]))) < tol:
                eigs.append(T[n - 1][n - 1])
                T = leading(T)
                n -= 1
                break
            mu = wilkinson(T)
            for i in range(n):
                T[i][i] -= mu
            Q, R = QR(T)
            T = zeros(n)
            for i in range(n):
                for j in range(max(0, i - 1), min(n, i + 2)):
                    for k in range(i, min(n, j + 2)):
                        T[i][j] += R[i][k] * Q[k][j]
            for i in range(n):
                T[i][i] += mu
    eigs.append(T[0][0])
    return sorted(eigs, reverse=True)
\end{lstlisting}
\end{document}